% =============================================================================
% Kapitel 4: Methodik
% =============================================================================

\section{Methodik}
\label{sec:methodik}

\subsection{Untersuchungsansatz}
\label{sec:methodik-ansatz}

Die Untersuchung folgt dem Prinzip der \textit{Fault Injection}: dem gezielten Einbringen
definierter Fehlerzustände in ein laufendes System, um dessen Reaktion empirisch zu
beobachten \cite{hsueh1997}. Für jeden Failure Case wird eine Hypothese über das erwartete
Systemverhalten formuliert, ein reproduzierbares Fehlerszenario hergestellt und das
tatsächliche Verhalten anhand definierter Messgrößen erfasst. Wo ein Fehlerzustand ohne
Eingriff in den Quellcode herstellbar ist, werden ausschließlich POSIX-Signale oder
externe Werkzeuge eingesetzt; wo ein Eingriff nötig ist, wird dieser minimal gehalten und
nach der Messung automatisch zurückgenommen.

\subsection{Verfahren der Fehlerinjektion}
\label{sec:methodik-injektion}

Jeder Failure Case bildet einen realen Fehlerzustand durch ein möglichst direktes
technisches Äquivalent nach. Tabelle~\ref{tab:injektion} ordnet die eingesetzten Verfahren
dem jeweils emulierten Fehler zu. Die ersten drei Verfahren kommen ohne jede
Code-Änderung aus: Ein \texttt{SIGKILL} beendet einen Prozess hart und entspricht damit
einem Absturz; \texttt{SIGSTOP} friert einen Prozess ein und bildet einen hängenden
Aufruf nach, etwa einen blockierenden Datenbank-Call; \texttt{pg\_terminate\_backend}
trennt eine Datenbankverbindung serverseitig. Nur zwei Cases erfordern einen minimalen,
nach der Messung automatisch zurückgenommenen Eingriff: das gezielte Auslösen eines
\texttt{OSError(28)} für die volle RAM-Disk sowie das Verschärfen von Timing-Parametern
zur Provokation der vermuteten Race Condition.

\begin{table}[htbp]
    \centering
    \caption{Injektionsverfahren und emulierte Fehlerzustände}
    \label{tab:injektion}
    \begin{tabularx}{\textwidth}{l>{\raggedright\arraybackslash}Xl}
        \hline
        \textbf{Verfahren} & \textbf{Emulierter Fehlerzustand} & \textbf{Case} \\
        \hline
        \texttt{SIGKILL}               & harter Prozessabsturz                         & 1 \\
        \texttt{OSError(28)}           & volle RAM-Disk (\texttt{No space left})       & 2 \\
        \texttt{SIGSTOP}/\texttt{SIGCONT} & eingefrorener Prozess (hängender DB-Call)  & 3 \\
        \texttt{pg\_terminate\_backend}   & serverseitig getrennte DB-Verbindung       & 4 \\
        Timing-Parameter                  & verschärftes Freigabe-Timing               & 5 \\
        \hline
    \end{tabularx}
\end{table}

\subsection{Testumgebung}
\label{sec:methodik-umgebung}

Alle Messungen erfolgen am Referenzsystem im Labor Bad Hersfeld. Das System wird über
einen Simulator betrieben, der die Kamerahardware durch das Einspielen vorab erfasster
Bilddaten ersetzt. Sämtliche Failure Cases sind dadurch ohne echte Kamerahardware unter
kontrollierten Bedingungen reproduzierbar. Tabelle~\ref{tab:testumgebung} fasst die
relevanten Eigenschaften zusammen.

\begin{table}[htbp]
    \centering
    \caption{Eigenschaften der Testumgebung}
    \label{tab:testumgebung}
    \begin{tabularx}{\textwidth}{lX}
        \hline
        \textbf{Eigenschaft} & \textbf{Wert} \\
        \hline
        Betriebssystem      & Linux 6.17 (x86\_64) \\
        Arbeitsspeicher     & 64 GB \\
        GPU                 & NVIDIA GeForce RTX 5090 (32 GB VRAM) \\
        Laufzeitumgebung    & Python 3.13 \\
        Startskript         & \texttt{main.py} (mit geänderter Konfiguration für Simulation) \\
        Bilddatenquelle     & 462 Bilder = 154 Inspektionen \\
        \hline
    \end{tabularx}
\end{table}

\subsection{Referenzmessung (Baseline)}
\label{sec:methodik-baseline}

Vor der Fehlerinjektion wurde eine Referenzmessung im Normalbetrieb ohne injizierte Fehler
durchgeführt. Sie dient als Vergleichsbasis für alle Failure Cases. Das System startet
neben dem Hauptprozess vier Service-Prozesse für LoggerProcess, FrameGrabber,
InspectionDataHandler und PipelineManager sowie einen Pool von zwölf Worker-Prozessen. Im
stationären Zustand erreicht es einen Durchsatz von
1,74~Inspektionen pro Sekunde bei einer mittleren Pipeline-Laufzeit von 383~ms. Die
zentrale Log-Queue bleibt dabei nahezu leer.

\subsection{Messgrößen}
\label{sec:methodik-messgroessen}

Für jeden Case werden vergleichbare Messgrößen erhoben:

\begin{itemize}
    \item \textit{Time-to-Detection (TTD):} Zeit von der Fehlerinjektion bis zur
          Erkennung durch das System (Shutdown oder Statusmeldung).
    \item \textit{Folgeverhalten:} qualitative Beschreibung der Reaktion der übrigen
          Prozesse (Weiterlauf, Deadlock, Leerschleife, Absturz).
    \item \textit{Queue-Tiefe und Speicherfüllstand über die Zeit:} quantitative
          Erfassung schleichender und kaskadierender Effekte.
    \item \textit{Graceful Degradation:} ob das System bei einem Teilausfall in einen
          kontrollierten Zustand übergeht oder vollständig ausfällt.
\end{itemize}

\subsection{Grenzen der Untersuchung}
\label{sec:methodik-grenzen}

Die Ergebnisse sind vor dem Hintergrund mehrerer methodischer Einschränkungen zu lesen.
Erstens erfolgen alle Messungen am Referenzsystem im Labor, nicht auf der späteren
Kundenhardware; die absoluten Zeiten sind daher nur eingeschränkt übertragbar, während die
zugrunde liegenden Fehlermechanismen von der Hardware unabhängig sind. Zweitens ersetzt der
Simulator die Kamerahardware: Hardwarespezifische Fehler wie ein Kamera-Disconnect werden
nicht direkt erzeugt, ihr Effekt auf das System aber über den Ausfall des FrameGrabbers
abgebildet. Drittens sind die Beobachtungsfenster mit 20 bis 120~Sekunden bewusst kurz
gehalten; schleichende Effekte über ganze Schichten werden aus den gemessenen Raten
hochgerechnet, nicht über die volle Dauer beobachtet. Die veraltete Datenbankverbindung
(Case~4) ist zudem aufgrund der maskierenden Entwicklungskonfiguration analytisch statt
empirisch belegt. Schließlich stützt sich jeder Case auf Einzelmessungen; der Fokus liegt
auf den qualitativen Fehlermechanismen und ihren Größenordnungen, nicht auf statistisch
abgesicherten Verteilungen.
